% =====================================================================
% Figura 1 — Relação entre os conceitos do referencial teórico.
%
% Substitui pipeline_referencial_sem_margem.png (981x560 px, 166 dpi).
% Desenhada em TikZ: saída vetorial, nítida em qualquer ampliação e
% editável diretamente no fonte.
%
% O painel 5 foi corrigido: a versão anterior listava o NNAPI entre os
% delegates avaliados, o que não corresponde ao escopo do trabalho.
%
% REQUER no preâmbulo (já presentes em artigo.tex):
%   \usepackage{tikz}
%   \usetikzlibrary{positioning, arrows.meta}
% =====================================================================

\begin{figure}[H]
    \centering
    \caption{Relação entre os principais conceitos abordados no referencial teórico}
    \label{fig:pipeline_referencial}

    \begin{tikzpicture}[
        painel/.style={
            rectangle, rounded corners=3pt, draw=black!40, line width=0.5pt,
            text width=3.75cm, inner sep=7pt, align=left,
            minimum height=4.0cm, anchor=north, font=\scriptsize
        },
        cabeca/.style={
            font=\scriptsize\bfseries, align=center,
            text width=4.80cm, anchor=south, inner sep=2pt,
            % hifenização desligada apenas nos títulos em caixa alta
            execute at begin node={\hyphenpenalty=10000 \exhyphenpenalty=10000 }
        },
        seta/.style={-{Stealth[length=5pt,width=3.5pt]},
                     draw=black!55, line width=0.7pt}
    ]

    \def\lx{5.00}   % passo horizontal entre colunas
    \def\ly{5.45}   % passo vertical entre linhas

    % ---------------------------------------------------------- linha 1
    \node[painel] (p1) at (0,0) {%
        Relevância econômica da cafeicultura brasileira.\par\smallskip
        Inspeção visual sujeita a variação de iluminação e a oclusão por folhas.\par\smallskip
        Detecção automática como apoio à decisão de colheita.\par};
    \node[cabeca, above=3pt of p1.north] {CAFEICULTURA E\\VISÃO COMPUTACIONAL};

    \node[painel] (p2) at (\lx,0) {%
        Detecção em estágio único, adequada a tempo real.\par\smallskip
        YOLOv8n, variante compacta para ambientes restritos.\par\smallskip
        YOLOv10n, arquitetura sem supressão não máxima.\par};
    \node[cabeca, above=3pt of p2.north] {DETECÇÃO DE OBJETOS\\COM YOLO};

    \node[painel] (p3) at (2*\lx,0) {%
        Execução local, sem dependência de conectividade.\par\smallskip
        Restrições de processamento, memória e energia.\par\smallskip
        Exigência de modelos de custo computacional reduzido.\par};
    \node[cabeca, above=3pt of p3.north] {INFERÊNCIA EM\\DISPOSITIVOS MÓVEIS};

    % ---------------------------------------------------------- linha 2
    \node[painel] (p4) at (0,-\ly) {%
        Formato otimizado para inferência embarcada.\par\smallskip
        Quantização pós-treinamento em FP16 e INT8.\par\smallskip
        Redução do tamanho do arquivo e da latência.\par};
    \node[cabeca, above=3pt of p4.north] {TENSORFLOW LITE};

    \node[painel] (p5) at (\lx,-\ly) {%
        Avaliação restrita aos \textit{delegates} de CPU e GPU.\par\smallskip
        NNAPI depreciada a partir do Android~15.\par\smallskip
        O suporte a operações varia conforme o acelerador.\par};
    \node[cabeca, above=3pt of p5.north] {\textit{DELEGATES}\\E ACELERAÇÃO};

    \node[painel] (p6) at (2*\lx,-\ly) {%
        Latência de inferência e latência total.\par\smallskip
        Estabilidade pelo coeficiente de variação.\par\smallskip
        Memória, consumo energético e tamanho do modelo.\par};
    \node[cabeca, above=3pt of p6.north] {MÉTRICAS DE AVALIAÇÃO};

    % ---------------------------------------------------------- ligações
    % altura fixa, para que as setas não se inclinem caso os painéis
    % tenham alturas ligeiramente diferentes
    \draw[seta] ([yshift=-2.0cm]p1.north east) -- ([yshift=-2.0cm]p2.north west);
    \draw[seta] ([yshift=-2.0cm]p2.north east) -- ([yshift=-2.0cm]p3.north west);
    \draw[seta] ([yshift=-2.0cm]p4.north east) -- ([yshift=-2.0cm]p5.north west);
    \draw[seta] ([yshift=-2.0cm]p5.north east) -- ([yshift=-2.0cm]p6.north west);

    \end{tikzpicture}

    \legend{Elaborada pela autora (2026).}
\end{figure}
